Thus we may assume additionally that
$W_2=W$. 
Denote $X=\Delta'[W]=\Delta'[W_2]$ and let $d\in W\cap D$.
By~\eqref{eq:Delta'} $\text{link}_d X$ is an induced subcomplex of $\Delta[C]$, and further, the suspension of $\text{link}_d X$ by the vertices $a$ and $b$ is an induced
subcomplex of $\Delta$.

Thus, if
$\widetilde{H}_{l-1}(\text{link}_d X)\neq 0$
then
$0\neq \widetilde{H}_l (\{\{a\},\{b\},\emptyset\}*\text{link}_d X)=\widetilde{H}_l(\Delta[\{a,b\}\cup V(\text{link}_d X)])$, and thus $\reg(I(G))\ge l+2$ as claimed. 
So we may assume $\widetilde{H}_{l-1}(\text{link}_d X)=0$.


Consider the Mayer--Vietoris long exact sequence
corresponding to the union
$X=\text{ast}_d X\cup_{\text{link}_d X} \bar{\text{St}}_d X$:
$$\widetilde{H}_l (\text{ast}_{d} X) \bigoplus \widetilde{H}_l (\overline{\text{St}}_{d} X) \longrightarrow \widetilde{H}_l (X) \longrightarrow \widetilde{H}_{l-1} (\text{link}_d X).$$
By assumption, the right term is zero while the middle term is nonzero, and closed stars have vanishing homology, hence
$\widetilde{H}_l (\text{ast}_{d} X)=
\widetilde{H}_l (\Delta'[W\setminus\{d\}])\neq 0$.
This contradicts the minimality of $W$, so in fact this last case cannot occur.
\end{proof}
